\documentclass[
    TitleLength=1,% 选择标题行数，范围是1~3
    Degree=UG,% 学位 本科：UG
    Type=thesis% thesis：毕业论文，design:毕业设计
]{njuthesis}

\addbibresource{njuthesis.bib}

\begin{document}

\maketitle

\begin{abstract}
    \zhlipsum[76][name=zhufu]
    \keywords{驾驶；星际穿越}
\end{abstract}

\begin{englishabstract}
    \blindtext
    % 英文关键词。关键词之间用英文半角逗号隔开，末尾无符号。
    \englishkeywords{Drive, Interstellar travel}
\end{englishabstract}

\tableofcontents

% 开始正文部分
\mainmatter

\chapter{绪论}
\section{简介}
南京大学学士学位论文\LaTeX 模板\cite{njuthesis}基于本科生院的论文撰写规范修改，用于生成符合南京大学学位论文排版要求和相应的国家规范、行业标准的学位论文。

此模板旨在为同学提供毕业论文书写的方便，如有模板问题或者版权问题，请于Github上提出issue。

\section{表格示例}
\begin{table}[htbp]
    \caption{经过测试的环境}
    \label{tab:mytab}
    \begin{tabular}{ccc}
        \toprule
        OS & Tex & 测试情况 \\
        \midrule
        Windows 10 & TexLive 2021 & ✔ \\
        Windows 10 & MikTex & ✔ \\
        Windows 10 & TexLive 2020 & ✔  \\
        Ubuntu 20.04 & TexLive 2021 & ✔ \\
        南大Tex & Overleaf & ✔ \\
        \bottomrule
    \end{tabular}
\end{table}
你可以使用\lstinline|table|环境插入标准三线表，如\cref{tab:mytab}所示，代码如下:
\begin{lstlisting}[language=TeX]
\begin{table}[htbp]
    \caption{经过测试的环境}
    \begin{tabular}{ccc}
        \toprule
        OS & Tex & 测试情况 \\
        \midrule
        Windows 10 & TexLive 2021 & ✔ \\
        Windows 10 & MikTex & ✔ \\
        Windows 10 & TexLive 2020 & ✔  \\
        Ubuntu 20.04 & TexLive 2021 & ✔ \\
        南大Tex & Overleaf & ✔ \\
        \bottomrule
    \end{tabular}
\end{table}}
\end{lstlisting}

\section{图片示例}
\begin{figure}[htbp]
    \includegraphics[width=0.5\textwidth]{njuname}
    \caption{南京大学名称}
    \label{fig:njuname}
\end{figure}
你可以使用\lstinline|figure|环境插入图片，如\cref{fig:njuname}，代码如下:
\begin{lstlisting}[language=TeX]
\begin{figure}
    \includegraphics[width=0.5\textwidth]{njuname}
    \caption{南京大学名称}
\end{figure}
\end{lstlisting}

\subsection{文字环绕图像}
% wrapfigure后面不能有空行
\begin{wrapfigure}{r}{0cm}
    % \label{fig:njulogo}
    \includegraphics[width=.3\textwidth]{njulogo}
    \caption{环绕图像测试}
\end{wrapfigure}
\zhlipsum[3][name=xiangyu]

\section{公式示例}
\begin{equation}\label{eq:dewitt}
    \int  e^{ax} \tanh bx\ dx = 
    \begin{cases}
    \displaystyle{ \frac{ e^{(a+2b)x}}{(a+2b)} 
    {_2F_1}\left[ 1+\frac{a}{2b},1,2+\frac{a}{2b}, -e^{2bx}\right] }& \\
    \displaystyle{
    \hspace{1cm}-\frac{1}{a}e^{ax}{_2F_1}\left[ 1, \frac{a}{2b},1+\frac{a}{2b}, -e^{2bx}\right]
    }
     & a\ne b \\
    \displaystyle{\frac{e^{ax}-2\tan^{-1}[e^{ax}]}{a} } & a = b
    \end{cases}
\end{equation}
    
你可以使用\lstinline|equation|环境插入公式，如\cref{eq:dewitt}，代码如下:
\begin{lstlisting}[language=TeX]
\begin{equation}\label{eq:dewitt}
    \int  e^{ax} \tanh bx\ dx = 
    \begin{cases}
    \displaystyle{ \frac{ e^{(a+2b)x}}{(a+2b)} 
    {_2F_1}\left[ 1+\frac{a}{2b},1,2+\frac{a}{2b}, -e^{2bx}\right] }& \\
    \displaystyle{
    \hspace{1cm}-\frac{1}{a}e^{ax}{_2F_1}\left[ 1, \frac{a}{2b},1+\frac{a}{2b}, -e^{2bx}\right]
    }
        & a\ne b \\
    \displaystyle{\frac{e^{ax}-2\tan^{-1}[e^{ax}]}{a} } & a = b
    \end{cases}
\end{equation}
\end{lstlisting}

% \chapter*{致谢}
% \addcontentsline{toc}{chapter}{致谢}
% 感谢\href{https://git.nju.edu.cn/nju-lug/lug-introduction}{NJU LUG}。

\printbibliography[heading=bibintoc, title=参考文献] % 方括号内的命令可以将参考文献列表加入目录
% \printindex

% 开始附录部分
\appendix 

\end{document}